\chapter{Penutup}
\section{Kesimpulan}
Berdasarkan analisis hasil simulasi dan evaluasi secara komprehensif yang telah dilakukan pada bab sebelumnya, penelitian ini menghasilkan tiga kesimpulan utama yang menjawab rumusan masalah sebagai berikut:
\begin{enumerate}
    \item Formulasi model Markowitz ke dalam kerangka EPG telah berhasil memetakan utilitas ekonomi menjadi sistem Hamiltonian fisis secara akurat dan konsisten. Penggunaan imbal hasi logaritmik untuk pembentukan matriks kovarians terbukti efektif menjaga simetri interaksi risiko ($J_{ij}$) tanpa distorsi, sementara imbal hasil sederhana mempertahankan linieritas bias ekspektasi imbal hasil ($h_i$). Implementasi suku penalti anggaran ($A$) dapat mendominasi koefisien kopling sehingga menciptakan batasan energi untuk mencegah sirkuit kuantum terjebak pada status portofolio yang melanggar aturan diversifikasi investasi.
    \item Penerapan ekuilibrium Nash sebagai inisialisasi awal mampu meningkatkan efisiensi dan stabilitas konvergensi algoritma VQE. Injeksi solusi ekuilibrium nash pada rotasi parameter awal memandu sirkuit melewati area dataran tandus yang dicirikan oleh kelenyapan varians gradien. Keunggulan strategis ini memungkinkan algoritma mencapai keadaan dasar pada kedalaman sirkuit yang dangkal ($L=2$ hingga $L=3$). Hal tersebut dapat mereduksi beban komputasi dan memicu keruntuhan entropi Von Neumann menuju konfigurasi probabilitas tunggal yang paling optimal.
    \item Arsitektur HEA yang dikombinasikan dengan algoritma SPSA menunjukkan performa komputasi yang efisien dalam mengeksekusi optimasi portofolio. Estimasi gradien menggunakan distribusi Rademacher memungkinkan sistem merespons volatilitas pasar melalui mekanisme penyeimbangan kembali. Hal ini divalidasi oleh hasil pengujian historis yang mana algoritma GT-VQE menghasilkan pertumbuhan modal lebih superior dibandingkan strategi klasik maupun VQE murni dengan nilai rasio Sharpe sebesar 1,0107 pada sistem $N=4$.
\end{enumerate}

\section{Saran}
Berdasarkan keterbatasan dan temuan dalam pengerjaan tugas akhir ini, terdapat beberapa saran untuk penelitian lanjutan di masa mendatang
\begin{enumerate}
    \item Mempertimbangkan parameter dearu sirkuit, seperti dekoherensi qubit dan fidelitas gerbang dalam mengeksekusi algoritma ini agar dapat diaplikasikan pada infrastruktur perangkat keras kuantum sesungguhnya.
    \item Menggunakan variasi struktur ansatz lain (seperti \textit{Qubit-Efficient} atau \textit{Variational Quantum Circuit}) serta membandingkan performa berbagai algoritma optimasi klasik lainnya untuk menemukan kombinasi hibrida yang memiliki laju konvergensi paling efisien dalam menangkap kompleksitas interaksi agen di pasar modal.
    \item Menggunakan rentang periode data historis yang lebih lama. Selain itu, pengujian pada periode yang sedang mengalami fenomena angsa hitam seperti dengan krisis ekonomi supaya dapat mengevaluasi agilitas dan reliabilitas algoritma dalam mempertahankan stabilitas ekuilibrium portofolio di tengah volatilitas pasar yang ekstrem.
\end{enumerate}
